\documentclass[12pt,letterpaper,oneside]{article}
\newcommand{\planCaso}{AM Taller Autocentro}
\newcommand{\planDescripcion}{Borrador académico basado en documentación local; datos comerciales por confirmar.}
\newcommand{\planFecha}{2026-09-20}
\newcommand{\planLugar}{Cajeme, Sonora; no acredita domicilio comercial}
\input{ITESCA/maestria-en-gestion-administrativa/plan-de-negocios-mga/formato-itesca-plan-negocios-generadas.tex}
\begin{document}
\portadaITESCA{Área geográfica de impacto del proyecto}{6547}
\section{Área geográfica de impacto del proyecto}

El presente documento constituye un \textbf{borrador hipotético pendiente de confirmación}. Su propósito es delimitar el procedimiento para determinar el área geográfica de impacto, estimar el mercado potencial y definir el mercado meta de AM Taller Autocentro. Debido a que todavía no se proporciona la localidad donde operaría el proyecto ni se ha seleccionado una fuente estadística concreta, no se atribuyen al negocio cifras de población, demanda, clientes, ventas, instalaciones o capacidad operativa comprobada.

AM Taller Autocentro se plantea como un proyecto local que integraría servicios de taller mecánico, llantera y refaccionaria, apoyados por una agenda de citas, control de inventario y revisión de proveedores. Por la naturaleza presencial de estas actividades, su alcance dependería de la cercanía, accesibilidad vial, tiempo de traslado y capacidad de atención que posteriormente se confirmen. Por ello, el análisis distingue entre población territorial, mercado potencial y segmento efectivamente atendible.

\subsection{Delimitación territorial propuesta}

La unidad geográfica principal está \textbf{pendiente de definir}. Antes de realizar cálculos deberán confirmarse el estado, municipio, localidad y, cuando corresponda, las colonias o códigos postales comprendidos. No se adopta Cajeme ni otra localidad como territorio obligatorio, ya que la información disponible para la actividad 6547 no establece una ubicación específica.

Una vez confirmada la ubicación de referencia, el área de impacto podría organizarse en tres niveles:

\begin{itemize}
    \item \textbf{Zona primaria:} colonias o sectores cercanos desde los cuales sería razonable acudir habitualmente al establecimiento propuesto. Sus límites deberán justificarse mediante vialidades, códigos postales, distancias o tiempos de traslado verificables.
    \item \textbf{Zona secundaria:} sectores cuyos residentes podrían desplazarse por conveniencia, disponibilidad de servicios o búsqueda de determinadas refacciones, aunque el recorrido sea mayor.
    \item \textbf{Zona complementaria:} localidades próximas cuya atención sería ocasional y dependería de la accesibilidad, capacidad operativa y conveniencia de las personas interesadas.
\end{itemize}

Esta clasificación deberá respaldarse con un mapa elaborado después de confirmar el punto de referencia del proyecto. La representación deberá identificar límites territoriales, vías principales, posibles barreras de acceso y un criterio uniforme de desplazamiento. El emblema circular dorado, la tipografía oscura y el lema ``Pasión por los autos. Compromiso contigo.'' podrán conservarse como elementos de identidad institucional del proyecto, sin modificar ni atribuir titularidad sobre la marca.

El impacto directo correspondería a las personas que pudieran contratar servicios de mecánica, llantas o refacciones. Los hogares, acompañantes y actividades relacionadas con la movilidad podrían formar parte del contexto indirecto, pero no deben contabilizarse como compradores independientes sin evidencia que lo justifique.

\subsection{Estimación del mercado potencial}

La investigación de mercados requiere recopilar información pertinente antes de seleccionar un segmento objetivo \autocite{itescaCompetencia2}. Para AM Taller Autocentro se propone combinar datos demográficos, indicadores vehiculares y evidencia sobre disposición y capacidad de compra mediante el siguiente procedimiento:

\begin{enumerate}
    \item Obtener la población total del territorio confirmado y registrar la institución responsable, nombre exacto del conjunto de datos, año de referencia, fecha de consulta y nivel geográfico.
    \item Identificar, cuando exista información oficial compatible, el número de hogares, población adulta, vehículos registrados u otro indicador relacionado con posibles usuarios.
    \item Establecer criterios de elegibilidad, como residencia o actividad habitual en la zona, uso o responsabilidad sobre un vehículo y posibilidad de trasladarse al área propuesta.
    \item Investigar la disposición de compra mediante una encuesta posterior sobre hábitos de mantenimiento, necesidades de servicio, criterios de selección y disposición para solicitar una cotización.
    \item Aproximar la capacidad de compra mediante indicadores estadísticos pertinentes y preguntas formuladas por rangos, sin recopilar información personal innecesaria.
    \item Excluir duplicidades y documentar las limitaciones de cada fuente antes de calcular el mercado.
\end{enumerate}

La fuente estadística concreta permanece pendiente. Podrían revisarse conjuntos oficiales del INEGI relacionados con población, hogares, unidades económicas o vehículos registrados, pero deberá seleccionarse la publicación correspondiente al territorio y periodo analizados. No deben mezclarse datos de años o niveles geográficos distintos sin explicar su compatibilidad y limitaciones.

La \textbf{población territorial} sería el total de residentes en el área delimitada. El \textbf{mercado potencial} comprendería únicamente a las personas que cumplan los criterios de elegibilidad, presenten una necesidad relacionada con los servicios propuestos y cuenten con disposición y capacidad de compra. Los vehículos no equivalen automáticamente a compradores: una persona puede utilizar varios vehículos y un hogar puede compartir uno. Mientras no se obtengan datos estadísticos y evidencia directa, la cantidad y proporción del mercado potencial deberán permanecer pendientes.

\subsection{Definición del mercado meta a cinco años}

El mercado meta será la parte del mercado potencial que AM Taller Autocentro podría atender con sus recursos efectivos. Su determinación requiere confirmar espacios de servicio, personal disponible, horarios, duración promedio de los trabajos, capacidad de inventario, tiempos de abastecimiento y presupuesto. Estos elementos son dependencias propuestas y no condiciones actualmente instaladas.

La proyección deberá formularse con los siguientes criterios:

\begin{itemize}
    \item \textbf{Año 1:} establecer una meta conservadora basada en la capacidad inicial confirmada.
    \item \textbf{Año 2:} revisar la utilización de espacios, recurrencia y tiempos efectivos de atención.
    \item \textbf{Año 3:} valorar una posible ampliación territorial o de servicios, únicamente si existe evidencia operativa.
    \item \textbf{Año 4:} analizar la consolidación del segmento atendido y los límites de capacidad.
    \item \textbf{Año 5:} definir una meta acumulada compatible con la infraestructura y los recursos autorizados.
\end{itemize}

En cada periodo deberán registrarse el mercado potencial expresado en personas, la meta anual de personas distintas y la proporción correspondiente. No se asignan cifras en este borrador porque una proyección sin población territorial ni capacidad documentada sería ficticia. La diferencia entre mercado disponible, mercado objetivo y capacidad de atención deberá conservarse de manera consistente con los conceptos empleados en el plan de negocios \autocite{planNegociosGlosario2024}.

\subsection{Control editorial y datos requeridos}

El resultado final deberá mantener trazabilidad con GGPN01 Plan de Negocios, de la Maestría en Gestión Administrativa del ITESCA, y con las especificaciones aplicables a la actividad 6547 \autocite{itescaPlanNegociosEspecificaciones2026}. Antes de su entrega deberán revisarse la correspondencia entre territorio y estadísticas, la vigencia de los datos, la ausencia de duplicidades y la coherencia entre mercado meta y capacidad operativa. También deberá incorporarse el mapa como evidencia gráfica, con fuente, año, límites y notas de elaboración claramente identificados.

Los datos necesarios para convertir este borrador en un trabajo real son: ubicación territorial exacta; colonias, localidades o códigos postales considerados; criterio de distancia o traslado; fuente estadística, edición y año; población de referencia; indicador vehicular pertinente; evidencia sobre disposición y capacidad de compra; recursos, horarios y capacidad operativa; así como formato, rúbrica y peso individual de la actividad. Sin la confirmación del territorio y de la fuente estadística no es posible calcular responsablemente el mercado potencial ni establecer metas numéricas para los próximos cinco años.
\clearpage
\printbibliography[title={Fuentes de consulta}]
\end{document}
